\chapter{User Requirements}\label{ch:requirements}

This chapter describes what PrepMeUp must do, seen from the people who use it. We
start with the users, then the user stories and the functional requirements, then
the requirements for quality, and finally the domain model, which shows the concepts
the requirements talk about. The full requirements specification is in Appendix 1.2.
Here we only show an overview and the requirements that matter most.

\section{Users and System Context}\label{sec:req-users}

PrepMeUp has two kinds of users. The \emph{customer} is a person who wants an
emergency stock at home. The customer signs up for a subscription in the customer
app, composes the orders and receives them. The \emph{administrator} is an employee of
PrepMeUp. The administrator maintains the package types the customers choose from and
sends the orders. \autoref{fig:rich-picture} in \autoref{ch:introduction} shows both
users and what moves between them and the system.

Both users log in through a separate identity provider. That is where they create an
account, log in and change their password, so PrepMeUp never handles a password. The
customer creates the account on the identity provider's login page and then enters
name and address in the app. We call that step onboarding. Only an existing
administrator can make another user an administrator. How the system and the identity
provider fit together is shown in Appendix 3.3.

Payment is simulated. The customer pays a subscription fee to a stored, fake card, and
no money is charged. There is no real grocery supplier either. Hardware, logistics and
the actual supply of goods are outside the project.

\section{Epics and User Stories}\label{sec:req-stories}

We have written the requirements as user stories, grouped into six epics.
\autoref{tab:epics} shows the epics and the functional requirements that realise them.

\begin{table}[H]
    \centering
    \begin{tabular}{@{}l p{5.6cm} l@{}}
        \hline
        Epic & What it covers & Requirements \\
        \hline
        Account & Creating an account, logging in, changing the password and deleting the account & FR-01, FR-08, FR-21 \\
        Order contents & Choosing the number of adult and child packages, their package types and extra equipment & FR-07, FR-09, FR-13 \\
        Subscription and orders & Getting a new order before the old one expires, seeing its status, pausing or cancelling & FR-01, FR-04, FR-07, FR-19 \\
        Notifications & Being told before the next order is sent & FR-20 \\
        Package catalogue & Adding package types and replacing them with new versions & FR-03, FR-11, FR-15 \\
        Dispatch & Seeing which orders to send and marking them as sent & FR-05, FR-06, NFR-10 \\
        \hline
    \end{tabular}
    \caption{The epics in the requirements specification and the requirements that realise them.}
    \label{tab:epics}
\end{table}

Three user stories describe the core of the product. The first is about the customer:

\begin{quote}
As a customer, I want to automatically receive a new order before the old one expires,
so that I don't have to keep track of the expiry date myself and have a chance to use
the soon-to-expire items.
\end{quote}

The two others are about the administrator:

\begin{quote}
As an administrator, I want to see the orders that should be sent, sorted by which order
should be sent first, so that orders are sent before the customer's previous order
expires.
\end{quote}

\begin{quote}
As an administrator, I want to mark an order as sent, so that the customer can see it is
on its way.
\end{quote}

Together they describe the loop that keeps the stock in date. The system plans the next
order, the administrator sends it, and the clock starts again.

\section{Functional Requirements}\label{sec:req-functional}

Each functional requirement has a number, FR-\emph{nn}, and a priority after the
MoSCoW method: must have, should have, could have or won't have. A requirement keeps
its number if it moves to another priority, so the numbers are not in order within
each priority. The must-have requirements are the smallest system that lets a customer
subscribe and receive a first order. They are shown in \autoref{tab:must-have}.

\begin{table}[H]
    \centering
    \begin{tabular}{@{}p{1.6cm} p{12.6cm}@{}}
        \hline
        Number & Requirement \\
        \hline
        FR-01 & The customer can create an account and sign up for a PrepMeUp subscription in the customer app. \\
        FR-02 & The back-end stores the customers, their subscriptions and their orders. \\
        FR-03 & There are default package types for a normal diet, one for adults and one for children. \\
        FR-04 & When an order is sent, its expiry date is set to one year after the sent date and kept for each of the customer's past orders. \\
        FR-05 & The administrator can see all planned orders, sorted by planned send date, earliest first. \\
        FR-06 & The administrator can register an order as sent. \\
        FR-19 & 14 days before an order's expiry date, the system creates the next order for a customer with an active subscription. It has the same package types and quantities as the previous order. \\
        \hline
    \end{tabular}
    \caption{The must-have functional requirements.}
    \label{tab:must-have}
\end{table}

The should-have requirements make the product more useful for the customer. The
customer can see the status of the orders and change a planned order (FR-07), change
name and address (FR-08), and choose which package types an order contains and how
many of each (FR-09). The customer also gets a notification when the next order is
created, so there is time to change it before it is sent (FR-20). The could-have
requirements include equipment packages such as a flashlight (FR-13) and a web build of
the customer app (FR-16).

Some requirements are won't have. A customer cannot choose how many days a package
lasts (FR-12). A customer who wants to cover more days orders more packages instead.
There are no packages made for allergies (FR-17) and no real grocery supplier (FR-18).
Deleting the account (FR-21) is also left out, because it is not needed to show the
core flow of the product.

The requirements changed along the way. The biggest change came in version 1.4, when
the orders went from being calculated from a household to being composed from package
types. We explain that change with the domain model in \autoref{sec:req-domain-model}.
In version 1.7 we set the expiry date to one year after the sent date, so it no longer
depends on the items in the order.

\section{Non-functional Requirements}\label{sec:req-nonfunctional}

The non-functional requirements are grouped by the categories in FURPS+. We have written
each of them so it can be measured and not just judged, and Appendix 6.1 describes how
each one is tested. \autoref{tab:nfr} gives a short version of each requirement.

\begin{table}[H]
    \centering
    \begin{tabular}{@{}p{2.6cm} p{1.6cm} p{10cm}@{}}
        \hline
        Category & Number & Requirement in short \\
        \hline
        Usability & NFR-01 & All text in the apps is in English, also error messages from the API. \\
                  & NFR-02 & An error message says what went wrong and what the user can do, without technical details. \\
        Reliability & NFR-03 & The customer app shows a connection error with an option to try again if the back-end cannot be reached. \\
                    & NFR-04 & The API answers 503 with an English message if the database cannot be reached. \\
        Performance & NFR-05 & With 10 users at the same time, 90\% of requests are answered in under 1 second. \\
        Portability & NFR-06 & The customer app runs as an Android app. \\
        Supportability & NFR-07 & Outside development, the API refuses to start without the address of the identity provider or the database. \\
                       & NFR-08 & Every API endpoint is covered by at least one automated test. \\
        Security & NFR-09 & A customer can only see and change their own data. \\
                 & NFR-10 & Endpoints for administrators reject customers. \\
                 & NFR-11 & Endpoints that need a logged-in user reject requests without a valid login. \\
        \hline
    \end{tabular}
    \caption{The non-functional requirements in short. The full wording is in Appendix 1.2.}
    \label{tab:nfr}
\end{table}

The security requirements matter most for the customer, because the system stores
their name and address. NFR-09 means that a customer can never see another customer's
orders, even by changing a number in a request.

\section{Domain Model}\label{sec:req-domain-model}

\autoref{fig:domain-model} shows the domain model. It shows the concepts the
requirements talk about and how they relate, without any technology. Appendix 2.2
describes it in more detail.

\drawiofig[\linewidth]{domain-model}{Domain model of PrepMeUp: the concepts from the requirements and how they are associated, without technology.}{fig:domain-model}

A customer lives at one address and has at most one subscription, and the subscription
schedules the orders. An order contains one or more packages. Each package is
described by a package type, for example an adult package or a child package, and the
administrator maintains the package types and sends the orders. A package type lists
its contents, which are items with a quantity. An item is either food or equipment, so
a flashlight is offered in the same way as food and water.

An order's status is not stored. It follows from whether the order has a sent date, so
an order is either planned or sent. The expiry date is stored on the order. It is one
year after the sent date, and the next order is created 14 days before it. The
customer has at most one payment method, and a notification tells the customer about
an order.

\subsection{From Household to Package Types}

In the first domain model, the customer described a household with members and a diet,
and the system calculated the package from it. In version 1.4 we removed the household,
and the customer now composes each order from package types instead.

We had already chosen packages with fixed contents, so the system did not need to know
who lives in the house, only how many packages of each type to send. That made the model
smaller. The diet became part of the package type, for example Vegetarian Adult.

The cost is that the customer now works out how many adult and child packages they
need. That is one of the barriers from \autoref{ch:introduction}, which the household
profile used to handle for them.
